% Fig. 2 block for paper/main.tex: the condensed (3x4, eight-panel) figure.
% Supersedes fig2_caption_4x4.tex on 2026-09-12.
%
% Letter map from the 4x4 caption: c->a, g+e->b, h+f->c, j+i->d, k->e,
% l->f, m->g, d+n->h. The (e) and (f) text is carried over from (k) and (l)
% unchanged apart from the letters. Conceptual level only: no resampling
% operator is named, per the 2026-08-09 decision.
%
% Two wording fixes folded into the rewrite because the sentences had to be
% rewritten anyway (both were cold-read conflicts between figure and caption):
%   - (a) drops the prime on cos(kr'), matching the legend;
%   - (d) states arctan(-Y/X) then "unwrapped", matching the panel title.
%
% This file is the working copy for further caption edits; re-apply the
% figure* block below to main.tex (after `git pull --rebase` in paper/) when
% it changes.

\begin{figure*}[tbp]
    \centering
    \includegraphics[width=\textwidth]{fig2_drift_correction_condensed.pdf}
    \caption{\textbf{Lawler--Fujita drift correction in a one-dimensional toy model.}
    (a) An ideal lattice signal $\cos(kr)$ with period $d=2$~nm (gray) and the measured signal $A(r)=\cos[k(r+u(r))]$ (orange), distorted by a nonlinear displacement field. In an experiment $u(r)$ is unknown and must be recovered from the signal itself; the toy uses $u(r)=\ln(r+1)/k$ so that the recovery in (d) can be checked against the true field. The measured signal is compressed where the drift changes rapidly; dots mark the lattice sites of each.
    (b,c) Multiplication by $\cos(kr)$ and $\sin(kr)$ (thin traces) separates the slowly varying drift phase $\theta(r)=k\,u(r)$ from fast $2k$ oscillations, and low-pass filtering isolates the slow quadrature components $X(r)$ and $Y(r)$ (thick), which follow $\tfrac{1}{2}\cos\theta$ and $-\tfrac{1}{2}\sin\theta$ (dashed).
    (d) The phase $\theta=\arctan(-Y/X)$, unwrapped, gives the recovered displacement field $u=\theta/k$ (dashed), shown against the true field (gray). Deviations within about 1~nm of the left edge in (b) to (d), and the flat segment near $r=0$ in (g), are boundary effects: the low-pass filter and the remap have no data beyond the edge, and the drift changes fastest there.
    (e) A single Fourier mode $\cos(qr)$, shown at the lattice wavevector $q=k$. The stems mark the mode's value at each sample's assigned position, i.e. the Fourier weight that sample carries in the sum for the coefficient $x_q$: a standard FFT pairs the measurement $y_n$ with the grid position $r_n$ (orange), while drift correction pairs the same $y_n$ with the corrected position $\tilde{r}_n = r_n + u(r_n)$ (blue), where the tilde marks a corrected quantity throughout; gray lines connect the two assignments of the same sample. The same measured value $y_n$ therefore multiplies the orange mode value in the uncorrected sum $x_q=\sum_n y_n\,e^{-iq r_n}$ and the blue mode value in the corrected sum $\tilde{x}_q=\sum_n y_n\,e^{-iq \tilde{r}_n}$, at the same frequencies $q$.
    (f) The corrected coefficients $\tilde{x}_q$ occupy the same uniform frequency grid as an ordinary FFT, so the standard inverse transform $F^{-1}$ evaluates the corrected map back on the uniform grid. Blue circles: the measured values $y_n$ registered at their corrected positions $\tilde{r}_n$ (in this single-mode illustration they trace the ideal lattice cosine). Open diamonds: the recomputed on-grid values $\tilde{y}_n = F^{-1}\tilde{x}$ evaluated at $r_n$; each is a sum over all samples rather than a moved copy of any single measurement, so no sample-to-sample connectors are drawn. In (e) and (f) only a subset of the full-grid sample positions is drawn for clarity.
    (g,h) Remapping the distorted signal (orange) by the recovered displacement restores the lattice periodicity (g, blue) and sharpens the Fourier peak (h); the ideal signal and its transform are shown in gray.}
    \label{fig:drift_correction_1d}
\end{figure*}
